\documentclass[10pt,a4paper]{amsart}
\usepackage[margin=19mm]{geometry}
\usepackage{amsmath,amssymb,mathtools}
\usepackage[scr=rsfs,cal=euler]{mathalfa}
\usepackage{xfrac}
\usepackage[unicode,hidelinks,bookmarks=false]{hyperref}
\newcommand{\ie}{i.e.\ }
\newcommand{\cf}{cf.\ }
\newcommand{\resp}{respectively\ }
\newcommand{\Subsection}{\subsection}
\providecommand{\nobreakdash}{\nobreak\hbox{-}}
\setlength{\emergencystretch}{2em}
\multlinegap0pt
\usepackage{graphicx}
\usepackage{tikz-cd}
% The author's own quiver style file (quiver.sty 1.4.2) is a thin wrapper around
% tikz-cd that also declares the TikZ styles his exported diagrams use and the two
% TikZ libraries they need; those declarations are carried into the wrapper verbatim.
\usetikzlibrary{calc}
\usetikzlibrary{decorations.pathmorphing}
\tikzset{curve/.style={settings={#1},to path={(\tikztostart)
    .. controls ($(\tikztostart)!\pv{pos}!(\tikztotarget)!\pv{height}!270:(\tikztotarget)$)
    and ($(\tikztostart)!1-\pv{pos}!(\tikztotarget)!\pv{height}!270:(\tikztotarget)$)
    .. (\tikztotarget)\tikztonodes}},
    settings/.code={\tikzset{quiver/.cd,#1}
        \def\pv##1{\pgfkeysvalueof{/tikz/quiver/##1}}},
    quiver/.cd,pos/.initial=0.35,height/.initial=0}
\tikzset{tail reversed/.code={\pgfsetarrowsstart{tikzcd to}}}
\tikzset{2tail/.code={\pgfsetarrowsstart{Implies[reversed]}}}
\tikzset{2tail reversed/.code={\pgfsetarrowsstart{Implies}}}
\tikzset{no body/.style={/tikz/dash pattern=on 0 off 1mm}}
\usepackage{mathrsfs}
\usepackage{stmaryrd}
% Class shim: the author writes the calligraphic letters of his own macros as
% {\cal X} with the article class, which keeps the old-style \cal switch; the AMS
% classes remove that switch, so it is restored here as an alias of \mathcal.
\providecommand{\cal}{\mathcal}
\allowdisplaybreaks

\makeatletter
\newcommand{\ostar}{\mathbin{\mathpalette\make@circled\star}}
\newcommand{\make@circled}[2]{%
  \ooalign{$\m@th#1\smallbigcirc{#1}$\cr\hidewidth$\m@th#1#2$\hidewidth\cr}%
}
\newcommand{\smallbigcirc}[1]{%
  \vcenter{\hbox{\scalebox{0.77778}{$\m@th#1\bigcirc$}}}%
}
\makeatother

\newtheorem{theorem}{Theorem}[section]
\newtheorem{lemma}[theorem]{Lemma}
\newtheorem{corollary}[theorem]{Corollary}
\newtheorem{proposition}[theorem]{Proposition}
\newtheorem{conjecture}[theorem]{Conjecture}

\theoremstyle{remark}
\newtheorem{remark}{Remark}[section]
\newenvironment{rmk}{\begin{remark}}{\hfill$\Diamond$\end{remark}}
\newtheorem{example}[remark]{Example}

\theoremstyle{definition}
\newtheorem{definition}[theorem]{Definition}

\DeclarePairedDelimiter\ceil{\lceil}{\rceil}
\DeclarePairedDelimiter\floor{\lfloor}{\rfloor}

\numberwithin{equation}{section}

\newcommand{\bvec}[1]{{\mathbf #1}}
\newcommand{\fder}[1]{\frac{\delta}{\delta j(#1)}}
\newcommand{\ind}[2]{\int d^{#1}\bvec{#2}}
\newcommand{\beq}{\begin{eqnarray}}
\newcommand{\eeq}{\end{eqnarray}}

\def\ker{\operatorname{ker}}
\newcommand{\ovo}{{^{(\bf{-1})}}}

\def\s{\mathfrak s}
\def\d{\mathfrak d}
\def\g{\mathfrak g}
\def\b{\mathfrak b}
\def\h{\mathfrak h}
\def\l{\mathfrak l}
\def\L{\mathscr{L}}
\def\i{\mathfrak i}
\def\so{\mathfrak{so}}
\def\an{\mathfrak{an}}
\def\su{\mathfrak{su}}
\def\G{\mathfrak{G}}
\def\D{\mathfrak{D}}
\def\t{\mathfrak{t}}

\def\bbZ{\mathbb{Z}}
\def\R{\mathbb{R}}
\def\bbC{\mathbb{C}}

\def\SU{\mathrm{SU}}
\def\SL{\mathrm{SL}}
\def\SO{\mathrm{SO}}
\def\AN{\mathrm{AN}}
\def\ISO{\mathrm{ISO}}
\def\rd{\mathrm{d}}

\def\demi{{\frac{1}{2}}}

\newcommand{\cA}{{\cal A}}
\newcommand{\cC}{{\cal C}}
\newcommand{\cD}{{\cal D}}
\newcommand{\cE}{{\cal E}}
\newcommand{\cF}{{\cal F}}
\newcommand{\cI}{{\cal I}}
\newcommand{\cT}{{\cal T}}
\newcommand{\cG}{{\cal G}}
\newcommand{\ctG}{{\tilde{\cal G}}}
\newcommand{\cB}{{\cal B}}
\newcommand{\cJ}{{\cal J}}
\newcommand{\cH}{{\cal H}}
\newcommand{\cM}{{\cal M}}
\newcommand{\cN}{{\cal N}}
\newcommand{\cO}{{\cal O}}
\newcommand{\cV}{{\cal V}}
\newcommand{\cS}{{\cal S}}
\newcommand{\cP}{{\cal P}}
\newcommand{\cZ}{{\cal Z}}
\newcommand{\cISO}{{\cal ISO}}

\newcommand{\sfx}{{\mathsf{x}}}
\newcommand{\sfy}{{\mathsf{y}}}

\makeatletter
\newsavebox{\@brx}
\newcommand{\llangle}[1][]{\savebox{\@brx}{\(\m@th{#1\langle}\)}%
  \mathopen{\copy\@brx\kern-0.5\wd\@brx\usebox{\@brx}}}
\newcommand{\rrangle}[1][]{\savebox{\@brx}{\(\m@th{#1\rangle}\)}%
  \mathclose{\copy\@brx\kern-0.5\wd\@brx\usebox{\@brx}}}
\makeatother

\makeatletter
\newcommand{\xRrightarrow}[2][]{\ext@arrow 0359\Rrightarrowfill@{#1}{#2}}
\newcommand{\Rrightarrowfill@}{\arrowfill@\equiv\equiv\Rrightarrow}
\newcommand{\xLleftarrow}[2][]{\ext@arrow 3095\Lleftarrowfill@{#1}{#2}}
\newcommand{\Lleftarrowfill@}{\arrowfill@\Lleftarrow\equiv\equiv}
\makeatother

\newcommand{\id}{\operatorname{id}}
\newcommand{\HC}[1]{\textcolor{red}{HC: #1}}

\setcounter{section}{6}
% Counter restoration: the very first Remark of this section, Remark 6.1, is a
% one-paragraph remark that stands on the section's opening page just above the
% delivered subsection and is therefore outside the unit; the remark counter is
% restored so that the remarks inside the unit carry the numbers of the version PDF.
\setcounter{remark}{1}
\begin{document}
\title[The belt-buckle moves for twist types on a fold]{Categorical quantum symmetries and ribbon tensor 2-categories}
\author{Hank Chen}
\date{Source: arXiv:2501.08041v2 (CC BY 4.0)}
\maketitle
\noindent\emph{Source and licence.} Categorical quantum symmetries and ribbon tensor 2-categories. Version arXiv:2501.08041v2 (CC BY 4.0). This material is distributed under the Creative Commons Attribution 4.0 International licence, http://creativecommons.org/licenses/by/4.0/, as recorded in the arXiv OAI licence metadata for this exact version and shown on the abstract page https://arxiv.org/abs/2501.08041v2. Full rights notice: licenses/arxiv-2501.08041-CC-BY-4.0.txt.\par\smallskip
\noindent\textit{Editorial scope.} Counted unit: original Proposition 6.3 of the article, the belt-buckle proposition (native environment \texttt{proposition}, no source label, source0.tex lines 2628--2634, printed as Proposition 6.3 between Proposition 6.1 and Corollary 6.2 of the same section), delivered together with its complete original author proof as the single primary proof body at source lines 2635--2684. The proof environment follows the statement immediately and no detached proof block is involved. No proof marker is added and no mathematics is written, repaired or re-derived. Necessary complete context is retained uncounted: source lines 2504--2627, that is the whole of the subsection ``Tortile objects; the ribbon balancing'' from its own subsection heading, with its first Reidemeister-move discussion, Proposition 6.1, the framing definition with its table and figures, Corollary 6.2, the two subsections on framed and unframed writhes, and the subsubsection ``Reidemeister II: double-twist cancellations and the belt-buckle move'' immediately preceding the counted statement. Nothing inside the delivered ranges is omitted except the twelve forward cross-references described below. The publisher's own class is not present in the runtime, so the article is delivered in the AMS amsart wrapper of the pinned TeX Live runtime together with the author's own theorem declarations, the author's own macro preamble and the author's own equation numbering scheme, with the section counter and the remark counter restored editorially so that the delivered subsection prints as 6.1 and the remarks it contains carry the numbers of the version PDF, the one-paragraph Remark 6.1 that the author places on the section's opening page just above the delivered subsection being outside the unit, and with the author's \texttt{\textbackslash numberwithin\{equation\}\{section\}} and per-section theorem counters reproduced unchanged, so that the counted statement prints as Proposition 6.3 and every numbered equation of the delivered range carries the number of the version PDF. The author's own \texttt{quiver} style file is a thin wrapper that loads \texttt{tikz-cd}, imports two TikZ libraries and declares the TikZ styles his exported diagrams use; the wrapper loads \texttt{tikz-cd}, imports those same libraries and carries those style declarations verbatim, so that every TikZ commutative diagram of the delivered ranges is delivered as native editable diagram code, not as an image. Class shim: the author writes the calligraphic letters inside his own macros, such as the object symbols, with the old-style \texttt{\textbackslash cal} font switch, which the article class keeps but the AMS classes remove; the wrapper restores that switch as an alias of \texttt{\textbackslash mathcal}, the very alphabet the author's articles render there, so that no macro body, symbol or letter of the author's source is altered. Macro faithfulness: every macro of the author's preamble span is carried into the wrapper verbatim except one, namely \texttt{\textbackslash insertpic}, a macro definition whose body embeds a graphical file inclusion and which therefore cannot appear in an editable flattened source; that macro is used nowhere in the article. Twelve cross-references of the delivered ranges point at results that lie outside the delivered unit, and they are delivered as link-only adaptations under a declared \texttt{reference\_only\_adaptation} receipt, each labelled with the number that the version PDF itself prints for that result: equations (2.1), (2.2), (3.4), (3.11), (3.12) and (5.4), Remark 3.4, Corollary 3.7, Warning 4.2 and subsection 5.2.1, together with the two occurrences of the subsection reference 5.2.1. The article's own bibliography is not delivered as source text: its \texttt{biblatex} data file is not reusable outside \texttt{biblatex}, so the five works that the delivered ranges cite are re-typeset from the reference list of the version PDF under the published numbers of that bibliography, namely 42, 44, 52, 85 and 116, so that the printed citation numbers coincide with the version PDF. Printed slips kept literal, in particular the misspelling ``quasi-Hemritian'' in the counted statement at source line 2629, the misspelling ``quasi-Hermimticity'' in the counted proof at source line 2683, and the misspelling ``Reidemseister'' at source line 2688. No figure, table, drawing or external asset is included as an image anywhere in the delivered ranges: the diagrams are editable TikZ code and the two tables of the subsection are native \texttt{tabular} material, so no artwork of any kind is redistributed and no bounded source comparison was needed.
\subsection{Tortile objects; the ribbon balancing}\label{ribbon}
We now finally come to the ribbon balancing/twist. Consider first the braiding structure $c_{\cD^*,\,^*\cD}: \cD^*\boxtimes\,^*\cD\rightarrow\,^*\cD\boxtimes\cD^*$. For the same reason as described in \S \href{https://arxiv.org/abs/2501.08041v2}{5.2.1}, these are associated with the following quantities
\begin{equation*}
    \mathfrak{t}= (-\cdot-)(\tilde{S}^{-1}\otimes\tilde{S})(\tilde R^T)=\frak{t}^l+\frak{t}^r,
\end{equation*}
which comes with its own nudging equations,
\begin{equation*}
    \frak{t}^l\cdot \frak{t}^r=\frak{t}^r\cdot \frak{t}^l
\end{equation*}
following from the compatibility of the cobraiding natural transformation $\tilde R$. We call these, specifically $\frak{t}$, the \textbf{tortile object} of $\tilde C$. Note $\tilde R$ being quasi-Hermitian does {\it not} imply $c_{\cD^*,\,^*\cD}^\dagger \cong c_{\,^*\cD,\cD^*}$, since the tortile object is different from the objects $\nu,\mu$ introduced in \S \href{https://arxiv.org/abs/2501.08041v2}{5.2.1}.

From the braiding map $c_{\cD^*,\,^*\cD}$, we define for each object $\cD$ the \textbf{left-over/right-under balancings}
\begin{align*}
    & \vartheta_\cD= (\overline{\operatorname{ev}}_\cD\boxtimes\cD^*)\circ (\cD\boxtimes c_{\cD^*,\,^*\cD})\circ(\operatorname{cev}_\cD\boxtimes\,^*\cD): \,^*\cD\rightarrow \cD^* \\
    & \bar\vartheta_\cD = (\,^*\cD\boxtimes{\operatorname{ev}}_\cD)\circ (c_{\cD^*,\,^*\cD}\boxtimes\cD)\circ(\cD^*\boxtimes\overline{\operatorname{cev}}_\cD):\cD^*\rightarrow\,^*\cD.
\end{align*}
The naturality of their composites gives, for each $\tilde C$-module functor $F:\cD\rightarrow\cD'$, the balancing 2-morphisms
% https://q.uiver.app/#q=WzAsNCxbMCwwLCJcXCxeKlxcY0QiXSxbMSwwLCJcXGNEXioiXSxbMCwxLCJcXCxeKlxcY0QiXSxbMSwxLCJcXGNEXioiXSxbMCwxLCJcXHZhcnRoZXRhX1xcY0QiXSxbMiwzLCJcXHZhcnRoZXRhX1xcY0QiLDJdLFswLDIsIlxcLF4qRiIsMl0sWzEsMywiRl4qIl0sWzIsMSwiXFx2YXJ0aGV0YV9GIiwyLHsic2hvcnRlbiI6eyJzb3VyY2UiOjIwLCJ0YXJnZXQiOjIwfSwibGV2ZWwiOjJ9XV0=
\[\begin{tikzcd}
	{\,^*\cD'} & {\cD'^*} \\
	{\,^*\cD} & {\cD^*}
	\arrow["{\vartheta_{\cD'}}", from=1-1, to=1-2]
	\arrow["{\,^*F}"', from=1-1, to=2-1]
	\arrow["{F^*}", from=1-2, to=2-2]
	\arrow["{\vartheta_F}"', shorten <=4pt, shorten >=4pt, Rightarrow, from=2-1, to=1-2]
	\arrow["{\vartheta_\cD}"', from=2-1, to=2-2]
\end{tikzcd},\qquad \begin{tikzcd}
	{\cD'^*} & {\,^*\cD'} \\
	{\cD^*} & {\,^*\cD}
	\arrow["{\bar\vartheta_{\cD'}}", from=1-1, to=1-2]
	\arrow["{F^*}"', from=1-1, to=2-1]
	\arrow["{\,^*F}", from=1-2, to=2-2]
	\arrow["{\bar\vartheta_F}"', shorten <=4pt, shorten >=4pt, Rightarrow, from=2-1, to=1-2]
	\arrow["{\bar\vartheta_\cD}"', from=2-1, to=2-2]
\end{tikzcd}.\]
One can show (through tedious but straightforward computations) from \href{https://arxiv.org/abs/2501.08041v2}{(3.11)} and \href{https://arxiv.org/abs/2501.08041v2}{(3.4)}, as well as \href{https://arxiv.org/abs/2501.08041v2}{(2.2)}, \href{https://arxiv.org/abs/2501.08041v2}{(2.1)}, that we have the following {balancing} equations
\begin{align}
    & \vartheta_{\cD\boxtimes\cA}\cong c_{\cD^*,\cA^*}\circ (\vartheta_\cD\boxtimes\vartheta_\cA)\circ c_{\,^*\cA,\,^*\cD},\nonumber\\
    & \bar\vartheta_{\cD\boxtimes\cA}\cong c_{\,^*\cD,\,^*\cA}\circ (\bar\vartheta_\cD\boxtimes\bar\vartheta_\cA)\circ c_{\cA^*,\cD^*}\label{balancing}
\end{align}
for each $\cD,\cA\in\operatorname{2Rep}(\tilde C;\tilde R)$; indeed, \href{https://arxiv.org/abs/2501.08041v2}{(3.4)} implies that these 2-morphisms are given by the unitary hexagonators, and so are unitary themselves and therefore invertible.

\begin{rmk}\label{pivottwist}
    Notice $\vartheta_{\cD^*}: \cD\rightarrow(\cD^*)^*$ is precisely the functor comparing $\cD$ and its double-dual mentioned in {\it Remark \href{https://arxiv.org/abs/2501.08041v2}{4.2}}. It is then clear that the 2-categorical dimension $\mathfrak{Dim}(\cD)$ admits an action only by the centralizer subcategory $C_{\operatorname{End}(\cD)}(\vartheta)$, which consist of functors $F\in\operatorname{End}(\cD)$ that commute with $\vartheta_{\cD^*}$, and natural transformations that commute with the 2-Drinfel'd modification $\omega_{\cD}$ (which we shall introduce soon). The issue raised in {\it Warning 2.2.5} of \cite{Douglas:2018} can thus be circumvented if $C_{\operatorname{End}(\cD)}(\vartheta)\simeq\mathsf{Hilb}$ is trivial. We call braided rigid tensor 2-categories with this property \textbf{maximally imbalanced}.
\end{rmk}


\subsubsection{2-Drinfel'd modifications}
Recall $\tilde S^2,\tilde S^{-2}:\tilde C\rightarrow\tilde C$ are by hypothesis monoidal autoequivalences. By combining \href{https://arxiv.org/abs/2501.08041v2}{(3.12)} and \textit{Remark \href{https://arxiv.org/abs/2501.08041v2}{3.4}}, we construct invertible 2-morphisms such that
\begin{equation}
    \overline{\operatorname{c/ev}}_\cD^{**}\Rightarrow \operatorname{c/ev}_{\cD^*},\qquad \,^{**}\operatorname{c/ev}_\cD\Rightarrow \overline{\operatorname{c/ev}}_{\,^*\cD},\label{dualinside}
\end{equation}
where "$\operatorname{c/ev}$" means either $\operatorname{ev}$ or $\operatorname{cev}$.
\begin{proposition}\label{2drinfeld}
     There are \textbf{2-Drinfel'd modifications}
        \begin{equation*}
        \omega_\cD: \bar\vartheta_\cD^*\Rightarrow \vartheta_{\cD^*},\qquad \bar\omega_\cD: \,^*\vartheta_{\cD}\Rightarrow \bar\vartheta_{\,^*\cD},
    \end{equation*}
    which witness the homotopy between left-over and right-under balancings upon a \textit{reversal} of the framing.
\end{proposition}
% \noindent 
\begin{proof}
Recall the (horizontal) antipode $\tilde S:\tilde C\rightarrow\tilde C^{\text{m-op},\text{c-op}}$ is a (strict) op-monoidal functor $\tilde S\circ (-\cdot -) = (-\cdot -)^\text{op}\circ (\tilde S\otimes\tilde S)$. Its adjunctions $\tilde S^{-1}\circ \tilde S\cong 1_{\tilde C} \cong \tilde S\circ\tilde S^{-1}$ lead to the following identifications
\begin{equation*}
    \tilde{S}\mathfrak{t}\cong (-\cdot -)(\tilde S^2\otimes 1)\tilde R,\qquad S^{-1}\mathfrak{t}\cong (-\cdot-)(1\otimes\tilde S^{-2})\tilde R
\end{equation*}
on the tortile objects. This in turn induces the invertible (unitary) 2-morphisms
\begin{equation}
    c_{-^*,-^*}: c_{\cD^*,\,^*\cD}^*\Rightarrow c_{(\cD^*)^*,\cD},\qquad c_{\,^*-,\,^*-}:\,^*c_{\cD^*,\,^*\cD}\Rightarrow c_{\cD,\,^*(\,^*\cD)},\label{dualbraid}
\end{equation}
corresponding to the braiding on the the duals $-^*$ or the pre-duals $\,^*-$. 

By horizontally composing $c_{-^*,-^*}$ with the 2-morphisms in {\bf Proposition \href{https://arxiv.org/abs/2501.08041v2}{3.7}}, we have with \eqref{dualinside},
\begin{align*}
    \bar\vartheta_\cD^* &= \big(\overline{\operatorname{cev}}^*_\cD\boxtimes (\cD^*)^*\big)\circ (\cD^*\boxtimes c_{\cD^*,\,^*\cD}^*) \circ (\operatorname{ev}_\cD^*\boxtimes\cD) \\
    &\Rightarrow \big(\operatorname{ev}_\cD\boxtimes(\cD^*)^*\big)\circ(\cD^*\boxtimes c_{(\cD^*)^*,\cD})\circ (\overline{\operatorname{cev}}_\cD^{**}\boxtimes\cD) \\
    &\cong \big(\overline{\operatorname{ev}}_{\cD^*}\boxtimes(\cD^*)^*\big)\circ (\cD^*\boxtimes c_{(\cD^*)^*,\cD})\circ (\operatorname{cev}_{\cD^*}\boxtimes\cD) = \vartheta_{\cD^*},
\end{align*}
where we have used the equivalences mentioned in {\it Remark \href{https://arxiv.org/abs/2501.08041v2}{3.4}}. This defines $\omega_\cD$. The other 2-morphism $\bar\omega_\cD$ can be constructed in an analogous way. 
\end{proof}
\noindent Moreover, the balancing \eqref{balancing} and \eqref{dualbraid} allow us to achieve the identifications
\begin{equation*}
    \omega_{\cD\boxtimes\cA} = c_{-^*,-^*}\circ (\omega_\cD\boxtimes\omega_\cA)\circ c_{-^*,-^*},\qquad \bar\omega_{\cD\boxtimes\cA} = c_{\,^*-,\,^*-}\circ (\bar\omega_\cD\boxtimes\bar\omega_\cA)\circ c_{\,^*-,\,^*-}
\end{equation*}
for each $\cD,\cA$, hence the 2-Drinfel'd modifications are compatible with the tensor product.

If we suppose for the moment that we have a monoidal pseudonatural isomorphism $-^{****}\simeq \id$ on $\operatorname{2Rep}(\tilde C;\tilde R)$, then the invertibility of \eqref{dualinside} gives rise to 
    \begin{equation*}
         \overline{\operatorname{c/ev}}_\cD \cong  (\overline{\operatorname{c/ev}})^{****}_\cD\cong \operatorname{c/ev}_{\cD^*}^{**},
    \end{equation*}
which in essence identifies barred quantities as a "half-way" to the quadruple dual. We will prove the triviality of the quadruple dual in a separate paper.

\begin{rmk}\label{categoricalS4}
   An alternative proof of $-^{****}\simeq \id$ is to leverage a "categorical" Radford $S^4$-formula\footnote{Recall that the usual Radford $S^4$-formula \cite{Delvaux2006ANO} states that the action of $S^4$ in a finite-dimensional Hopf algebra is an inner automorphism by grouplike elements.} proven in \cite{Etingof:2004}: there exists a \textit{distinguished invertible object} $c\in C$ in a finite tensor category $C$ with a natural isomorphism $\delta: -^{**} \Rightarrow c\otimes \,^{**}-\otimes c^{-1}$ --- the idea is that the rigid duality $-^*$ gives rise to an antipode functor on the dual $C^*$. Thus it is reasonable to expect a categorical Radford $S^4$-formula for the Hopf category $\tilde{C}=\tilde C$ to hold. One subtlety is that $\tilde C$ itself is not finite, but it can be made unimodular with the use of a \textit{Haar cointegral} (see \S 5.2.3 in \cite{chen2:2025}).
   % In analogy with the double construction of Majid \cite{Majid:1994nw,Majid:1996kd}, we may achieve this through the Frobenius pairing introduced in \cite{Crane:1994ty}.
\end{rmk}   

Due to the naturality of the ribbon balancing functors, an immediate corollary is the following.
\begin{corollary}\label{generalC8}
    When $\omega_\cD,\bar\omega_\cD$ are invertible, then for all $\tilde C$-module functors $F:\cD\rightarrow\cD'$ we have
    \begin{equation}
    \vartheta_{F^*} = \bar\vartheta_F^*,\qquad \vartheta_{\,^* F} = \,^*\bar\vartheta_F.\label{rigidC8}
\end{equation}
As such, $\operatorname{End}(\cI)$ is a ribbon tensor category.
\end{corollary}
\noindent Notice these are rigid dagger generalizations of the C8 condition in \cite{Douglas:2018}, Definition 2.2.4.

\begin{rmk}
    The reason we call $\omega_\cD,\bar\omega_\cD$ the "2-Drinfel'd modifications" is the following. In the scenario where $\vartheta,\bar\vartheta$ are genuine "twists" --- ie. (pseudo)natural transformations on the identity on $\operatorname{2Rep}(\tilde C;\tilde R)$ (which cannot happen unless the duality is involutive) --- then  $\omega=\bar\omega^{-1}$ is an invertible modification between these pseudonatural twists. As such, they serve to "change" the pivotal structures on rigid 2-categories, and hence appears in \textit{Remark \ref{pivottwist}} as part of the notion of "maximal imbalancing".
\end{rmk}


% \begin{rmk}
%     Note the conditions \eqref{adjdual1}, though more general than \eqref{adjdual}, is still geometrically 
% \end{rmk}

\subsubsection{Reidemeister II: double-twist cancellations and the belt-buckle move}\label{beltbuckle}
In the rigid dagger setting, each of the composites in the definition of the ribbon balancings $\vartheta_\cD,\bar\vartheta_\cD$ are planar-unitary, and hence are themselves planar-unitary. As such they admit folds against their adjoints. The folds
\begin{equation*}
    e_{\vartheta_\cD}:\vartheta_\cD^\dagger\circ\vartheta_\cD\Rightarrow 1_{\,^*\cD},\qquad e_{\bar\vartheta_\cD}: \bar\vartheta_\cD^\dagger\circ\bar\vartheta_\cD\Rightarrow 1_{\cD^*},
\end{equation*}
in particular, are known as the \textbf{Kauffman double twist cancellations}; see fig. 55 (d) in \cite{Barrett_2024}. These adjunctions witness the cancellation of twists that live on the {\it same side} of the tangle. We call such twist cancellations "of the first type". 

As the name suggests, there is another type of double twist cancellation, in which the twists lie on different sides of the tangle. 

\begin{proposition}
    Suppose $\tilde R$ is quasi-Hemritian \href{https://arxiv.org/abs/2501.08041v2}{(5.4)}. There are 2-morphisms
    \begin{equation*}
    \kappa_\cD:\big(\cD\boxtimes(\vartheta_\cD\circ\bar\vartheta_\cD)\big)\circ\operatorname{cev}_\cD\Rightarrow \operatorname{cev}_\cD,\qquad \bar\kappa_\cD:\big((\bar\vartheta_\cD\circ \vartheta_\cD)\boxtimes\cD\big)\circ\overline{\operatorname{cev}}_\cD\Rightarrow\overline{\operatorname{cev}}_\cD
    \end{equation*}
    which witness the null-homotopy of distinct twist types on a fold.
\end{proposition}

\begin{proof}
    To prove this, we need to introduce the so-called \textbf{belt-buckle moves}. Geometrically, these are isotopies which "drag" twists on the \textit{outside} (and only the outside!) of a fold to the top. They are implemented by the following 2-morphisms
\begin{equation*}
    \mathrm{P}_\cD: (\cD\boxtimes\bar\vartheta_\cD)\circ\operatorname{cev}_\cD\Rightarrow c_{\,^*\cD,\cD}\circ\overline{\operatorname{cev}}_{\cD},\qquad \mathrm{Q}_\cD: (\vartheta_{\cD}\boxtimes \cD)\circ \overline{\operatorname{cev}}_\cD\Rightarrow c_{\cD,\cD^*}\circ \operatorname{cev}_\cD.
\end{equation*}
We shall describe the construction of $\mathrm{P}_\cD$ in our current context; $\mathrm{Q}_\cD$ can be obtained in a similar manner. From the quasi-Hermitian hypothesis $c_{\,^*\cD,\cD}^\dagger\cong c_{\cD,\,^*\cD}$, we have an adjunction
% https://q.uiver.app/#q=WzAsMixbMCwwLCIoXFwsXipcXGNEKVxcY0RcXGNEXipcXGNEIl0sWzIsMCwiXFxjRChcXCxeKlxcY0QpXFxjRF4qXFxjRCJdLFswLDEsImNfe1xcLF4qXFxjRCxcXGNEfVxcY0ReKlxcY0QiLDIseyJjdXJ2ZSI6Mn1dLFsxLDAsImNfe1xcY0QsIFxcLF4qXFxjRH1cXGNEXipcXGNEIiwyLHsiY3VydmUiOjJ9XSxbMiwzLCJcXHBlcnAiLDIseyJzaG9ydGVuIjp7InNvdXJjZSI6MjAsInRhcmdldCI6MjB9fV1d
\[e_{c_{\,^*\cD,\cD}}\cD^*\cD: \begin{tikzcd}
	{(\,^*\cD)\cD\cD^*\cD} && {\cD(\,^*\cD)\cD^*\cD}
	\arrow[""{name=0, anchor=center, inner sep=0}, "{c_{\,^*\cD,\cD}\cD^*\cD}"', curve={height=12pt}, from=1-1, to=1-3]
	\arrow[""{name=1, anchor=center, inner sep=0}, "{c_{\cD, \,^*\cD}\cD^*\cD}"', curve={height=12pt}, from=1-3, to=1-1]
	\arrow["\vdash"{description}, draw=none, shorten <=3pt, shorten >=3pt, Rightarrow, from=0, to=1]
\end{tikzcd}: \iota_{c_{\,^*\cD,\cD}}\cD^*\cD.\]
The 2-morphism $\mathrm{P}_\cD$ is given by the diagram
% https://q.uiver.app/#q=WzAsNyxbMSwwLCJcXGNJIl0sWzAsMSwiXFxjRFxcY0ReKiJdLFsxLDIsIlxcY0RcXGNEXiooXFwsXipcXGNEKVxcY0QiXSxbMiwxLCIoXFwsXipcXGNEKVxcY0QiXSxbMiwyLCIoXFwsXipcXGNEKVxcY0RcXGNEXipcXGNEIl0sWzMsMSwiXFxjRChcXCxeKlxcY0QpIl0sWzMsMiwiXFxjRChcXCxeKlxcY0QpXFxjRF4qXFxjRCJdLFswLDNdLFszLDJdLFswLDFdLFsxLDJdLFswLDIsIlxcdXBzaWxvbl97XFxvdmVybGluZXtcXG9wZXJhdG9ybmFtZXtjZXZ9fV9cXGNELFxcb3BlcmF0b3JuYW1le2Nldn1fXFxjRH0iLDAseyJzaG9ydGVuIjp7InNvdXJjZSI6MjAsInRhcmdldCI6MjB9LCJsZXZlbCI6Mn1dLFs0LDNdLFszLDVdLFsyLDRdLFsyLDYsIiIsMix7ImN1cnZlIjo1fV0sWzYsNV0sWzQsNSwiXFx1cHNpbG9uX3tcXG9wZXJhdG9ybmFtZXtldn1fXFxjRCxjX3tcXCxeKlxcY0QsXFxjRH19IiwwLHsic2hvcnRlbiI6eyJzb3VyY2UiOjIwLCJ0YXJnZXQiOjIwfSwibGV2ZWwiOjJ9XSxbNCw2LCIiLDAseyJjdXJ2ZSI6LTF9XSxbNiw0LCIiLDAseyJjdXJ2ZSI6LTF9XSxbMTUsNCwiXFxPbWVnYV97XFwsXipcXGNEfFxcY0RcXGNEXip9IiwyLHsic2hvcnRlbiI6eyJzb3VyY2UiOjIwLCJ0YXJnZXQiOjEwfX1dLFszLDE0LCIiLDIseyJzaG9ydGVuIjp7InNvdXJjZSI6MjAsInRhcmdldCI6MjB9fV0sWzE4LDE5LCJcXGNvbmciLDEseyJzaG9ydGVuIjp7InNvdXJjZSI6MjAsInRhcmdldCI6MjB9LCJzdHlsZSI6eyJib2R5Ijp7Im5hbWUiOiJub25lIn0sImhlYWQiOnsibmFtZSI6Im5vbmUifX19XV0=
\[\mathrm{P}_\cD= \begin{tikzcd}
	& \cI \\
	{\cD\cD^*} && {(\,^*\cD)\cD} & {\cD(\,^*\cD)} \\
	& {\cD\cD^*(\,^*\cD)\cD} & {(\,^*\cD)\cD\cD^*\cD} & {\cD(\,^*\cD)\cD^*\cD}
	\arrow[from=1-2, to=2-1]
	\arrow[from=1-2, to=2-3]
	\arrow["{\upsilon_{\overline{\operatorname{cev}}_\cD,\operatorname{cev}_\cD}}", shorten <=6pt, shorten >=6pt, Rightarrow, from=1-2, to=3-2]
	\arrow[from=2-1, to=3-2]
	\arrow[from=2-3, to=2-4]
	\arrow[from=2-3, to=3-2]
	\arrow[""{name=0, anchor=center, inner sep=0}, from=3-2, to=3-3]
	\arrow[""{name=1, anchor=center, inner sep=0}, curve={height=30pt}, from=3-2, to=3-4]
	\arrow[from=3-3, to=2-3]
	\arrow[shorten <=8pt, shorten >=8pt, Rightarrow, from=3-3, to=2-4]
	\arrow[""{name=2, anchor=center, inner sep=0}, draw=none, curve={height=-6pt}, from=3-3, to=3-4]
	\arrow[from=3-4, to=2-4]
	\arrow[""{name=3, anchor=center, inner sep=0}, draw=none,curve={height=-6pt}, from=3-4, to=3-3]
	\arrow[shorten <=6pt, shorten >=6pt, Rightarrow, from=2-3, to=0]
	\arrow["{\Omega_{\,^*\cD|\cD\cD^*}}"', shorten <=3pt, shorten >=1pt, Rightarrow, from=1, to=3-3]
	\arrow["\vdash"{description}, draw=none, from=2, to=3]
\end{tikzcd},\]
where the central triangle is filled by $c_{\,^*\cD\boxtimes\cD,\operatorname{cev}_\cD}\bullet \bar K_\cD'$, and the square to its right is ${\upsilon_{\operatorname{ev}_\cD,c_{\,^*\cD,\cD}}}$. 

Now starting with the functor 
\begin{equation*}
    \big(\cD\boxtimes(\vartheta_\cD\circ\bar\vartheta_\cD)\big)\circ\operatorname{cev}_\cD: \cI \rightarrow \cD\boxtimes\,^*\cD,
\end{equation*}
we can apply a series of 2-morphisms
\begin{align*}
    \big(\cD\boxtimes(\vartheta_\cD\circ\bar\vartheta_\cD)\big)\circ\operatorname{cev}_\cD &\xRightarrow{(\cD\boxtimes\vartheta_\cD)\circ \mathrm{P}_\cD} (\cD\boxtimes\vartheta_\cD)\circ c_{\,^*\cD,\cD}\circ\overline{\operatorname{cev}}_\cD \\
    & \xRightarrow{c_{\cD,\vartheta_\cD}\circ\overline{\operatorname{cev}}_\cD} c_{\cD^*,\cD}\circ (\vartheta_\cD\boxtimes\cD)\circ\overline{\operatorname{cev}}_\cD\xRightarrow{c_{\cD^*,\cD}\circ \mathrm{Q}_\cD}c_{\cD^*,\cD}\circ c_{\cD,\cD^*}\circ\operatorname{cev}_\cD\\
    & \cong c_{\cD,\cD^*}^\dagger\circ c_{\cD,\cD^*}\circ\operatorname{cev}_\cD \xRightarrow{e_{c_{\cD,\cD^*}}\circ\overline{\operatorname{cev}}_\cD} \overline{\operatorname{cev}}_\cD,
\end{align*}
where in the last line we have used the quasi-Hermimticity property $c_{\cD^*,\cD} \cong c_{\cD,\cD^*}^\dagger$. This defines $\kappa_\cD$; the other one $\bar\kappa_\cD$ can be obtained similarly. 
\end{proof}

\begin{thebibliography}{99}
\bibitem[42]{Douglas:2018} Christopher L. Douglas and David J. Reutter. ``Fusion 2-categories and a state-sum invariant for 4-manifolds''. In: \emph{arXiv: Quantum Algebra} (2018). url: \url{https://api.semanticscholar.org/CorpusID:119305837}.
\bibitem[44]{Barrett_2024} John W. Barrett, Catherine Meusburger, and Gregor Schaumann. ``Gray categories with duals and their diagrams''. In: \emph{Advances in Mathematics} \textbf{450} (July 2024), p. 109740. issn: 0001-8708. doi: 10.1016/j.aim.2024.109740.
\bibitem[52]{Etingof:2004} P. Etingof et al. \emph{Tensor Categories}. Mathematical Surveys and Monographs. American Mathematical Society, 2016. isbn: 9781470434410.
\bibitem[85]{chen2:2025} Hank Chen. ``Combinatorial quantization of 4d 2-Chern-Simons theory II: Quantum invariants of higher ribbons in $D^4$''. In: (June 2025). arXiv: 2506.05785 [math-ph].
\bibitem[116]{Delvaux2006ANO} Lydia Delvaux, Alfons Van Daele, and Shuanhong Wang. ``A note on Radford's $S^4$ formula''. In: \emph{arXiv: Rings and Algebras} (2006). url: \url{https://api.semanticscholar.org/CorpusID:1770960}.
\end{thebibliography}
\end{document}
